% Source: front matter for the volume. Structure follows the Gerzensee course
% description in Fukuda (2026), "Mathematics Review Course," slides 2-6, and the
% "Structure of the Lecture Notes" section of Fukuda (2026), pp. xi-xiii.

\chapter*{Preface}
\addcontentsline{toc}{chapter}{Preface}
\markboth{Preface}{Preface}

A first-year economics doctoral student meets three kinds of mathematical
difficulty. The first is technical: an argument uses the implicit function
theorem, or the Banach fixed point theorem, or the fact that a continuous
function on a compact set attains its maximum, and the student has either never
seen the result or has seen it without its hypotheses. The second is
translational: the mathematics is familiar but the economic object it stands for
is not, so that a Lagrange multiplier is manipulated correctly without being
recognised as a shadow price and a fixed point is computed without being
recognised as an equilibrium. The third is architectural: the results are known
individually but not as a system, so that the student does not see that the
existence of a Walrasian equilibrium, the existence of a value function solving
a Bellman equation, and the existence of a least-squares projection are three
instances of the same handful of theorems about compactness, completeness, and
convexity.

These notes address all three. They follow the topics of the Mathematics Review
Course of the Swiss Program for Beginning Doctoral Students in Economics, taught
at the Study Center Gerzensee: probability and conditional expectation, calculus
in one and several variables, static optimisation, real analysis, and
discrete-time dynamic programming \autocite{fukuda2026logistics}. That course
compresses these subjects into four days, and its slides are deliberately
economical --- they state what is needed, illustrate it, and move on. A student
who wants to know why a hypothesis is there, what fails without it, or how a
result generalises must look elsewhere. The purpose of this volume is to be that
elsewhere, while keeping the course's sequence, its notation, and above all its
economic motivation.

The mathematical standard aimed at is that of a graduate text rather than that
of a review course: hypotheses are attached to the space on which a statement is
made, necessary conditions are distinguished from sufficient ones, existence from
uniqueness, and local statements from global ones. Where a teaching-level
simplification appears in the course slides, the simplification is retained for
its intuitive value and the exact statement is given alongside it, with the
relation between the two made explicit; the Kuhn--Tucker theorem in
\autoref{ch:optimization} and the envelope formula in \autoref{ch:dynamic} are
the clearest instances. A result that is quoted rather than proved is marked as
such, with the source and the reason for the omission given where it is used.

The economic applications are worked rather than named. Each one specifies the
environment, the preferences or technology, the constraint set, the
correspondence between mathematical and economic objects, the derivation, and
what goes wrong when a hypothesis fails.
The applications are drawn from the course itself: profit maximisation by
price-taking and price-setting firms, the risk premium and the Arrow--Pratt
coefficient, ordinary least squares as an orthogonal projection, the
consumer-choice problem, intertemporal consumption and the Euler equation, Nash
bargaining, Bayesian updating in a market-microstructure setting, normal-normal
learning, and the cake-eating, Ramsey growth, and Lucas asset-pricing problems.

\section*{Sources}

Three layers of source material stand behind the volume.

The first is the set of six slide decks of the Gerzensee course
\autocite{fukuda2026probability,fukuda2026calculus,fukuda2026optimization,%
fukuda2026analysis,fukuda2026dynamic,fukuda2026logistics}. These fix the scope:
the topics are the course's topics, and the results the slides state are carried
over here with their hypotheses supplied.

The second is the accompanying reference volume
\autocite{fukuda2026notes}, a 768-page set of lecture notes of which the slides
are a distillation. It supplies definitions the slides state informally, proofs
the slides sketch, exercises the slides mark optional, and the technical
conditions the slides suppress. Where a result appears in both, the numbering of
the reference volume is given so that the two can be read together.

The third consists of the mathematical texts that carry each topic beyond
what a review course can reach:
\textcite{rudin1976principles} and \textcite{amann2005analysisI,%
amann2008analysisII,amann2009analysisIII} for real analysis;
\textcite{axler2024linear} for linear algebra;
\textcite{boyd2004convex} and \textcite{rockafellar1970convex} for convexity and
optimisation; \textcite{legall2022measure} and \textcite{legall2016brownian} for
measure-theoretic probability; \textcite{corbae2009introduction} and
\textcite{ok2007real} for analysis written with economic applications in view;
and \textcite{ljungqvist2018recursive} and \textcite{stokey1989recursive} for
recursive methods. Each is cited, with the relevant chapter or section, at the
points where its content is used. \textcite{legall2016brownian} is listed
because it is the natural continuation for readers who go on to
continuous-time finance; nothing in this volume requires stochastic calculus,
and it is therefore not cited again.

\chapter*{How to Use These Notes}
\addcontentsline{toc}{chapter}{How to Use These Notes}
\markboth{How to Use These Notes}{How to Use These Notes}

\section*{Order of reading}

The chapters are ordered by logical dependence, not by the order in which the
four-day course presents them. The course opens with probability because Bayes'
rule is immediately useful; but the sharpest statement about conditional
expectation --- that it is the orthogonal projection of a square-integrable
random variable onto the space of functions of the conditioning variable --- needs
completeness, convexity, and a projection theorem. Those are developed in
Chapters~\ref{ch:completeness} and~\ref{ch:convex}, so probability appears here as
\autoref{ch:probability}. The dependence structure is as follows.

\begin{itemize}
	\item \textbf{Foundations.} \autoref{ch:sets} (sets, relations, logic, the real
	      number system) and \autoref{ch:linear} (vector spaces, linear maps,
	      matrices, spectra) presuppose nothing.
	\item \textbf{Analysis.} \autoref{ch:metric} (inner products, norms, metrics)
	      requires \autoref{ch:linear}. \autoref{ch:topology} (open and closed sets,
	      sequences, compactness) requires \autoref{ch:metric}.
	      \autoref{ch:continuity} (continuity, semicontinuity, the Weierstrass
	      theorem) requires \autoref{ch:topology}. \autoref{ch:completeness}
	      (completeness, the contraction mapping theorem, Blackwell's conditions)
	      requires \autoref{ch:continuity}.
	\item \textbf{Calculus.} \autoref{ch:calculus1} (one variable) requires only
	      \autoref{ch:sets} and the limit notions of \autoref{ch:topology}.
	      \autoref{ch:calculusn} (several variables, the implicit function theorem)
	      requires \autoref{ch:calculus1}, \autoref{ch:linear} and, for the proof of
	      the implicit function theorem, \autoref{ch:completeness}.
	      \autoref{ch:integration} (integration) requires \autoref{ch:calculus1}.
	\item \textbf{Convexity and optimisation.} \autoref{ch:convex} (convex sets and
	      functions, separation) requires \autoref{ch:topology} and
	      \autoref{ch:calculusn}. \autoref{ch:projection} (projection, orthogonality,
	      least squares) requires \autoref{ch:completeness} and \autoref{ch:convex}.
	      \autoref{ch:optimization} (Lagrange and Kuhn--Tucker) requires
	      \autoref{ch:convex} and \autoref{ch:calculusn}. \autoref{ch:maximum}
	      (correspondences, the maximum theorem, fixed points) requires
	      \autoref{ch:continuity} and \autoref{ch:convex}.
	\item \textbf{Probability and dynamics.} \autoref{ch:probability} requires
	      \autoref{ch:integration} and \autoref{ch:projection}.
	      \autoref{ch:dynamic} (discrete-time dynamic programming) requires
	      \autoref{ch:completeness}, \autoref{ch:optimization}, \autoref{ch:maximum},
	      and, for the stochastic case, \autoref{ch:probability}.
	      \autoref{ch:continuoustime} (continuous-time optimisation) requires
	      \autoref{ch:calculusn} and \autoref{ch:dynamic}.
\end{itemize}

A reader following the four-day course can read
\autoref{ch:probability}, then Chapters~\ref{ch:calculus1},
\ref{ch:calculusn} and~\ref{ch:integration}, then \autoref{ch:optimization},
then Chapters~\ref{ch:metric}--\ref{ch:completeness}, then
\autoref{ch:dynamic}, consulting earlier chapters as the cross-references
demand.

\section*{Conventions used throughout}

Definitions are boxed in green, formal statements in blue, and remarks and
examples are set plainly. A statement labelled \textbf{Assumption} is a standing
hypothesis that remains in force until it is explicitly withdrawn; a statement
labelled \textbf{Theorem}, \textbf{Proposition}, \textbf{Lemma}, or
\textbf{Corollary} carries all of its hypotheses within it and depends on no
standing assumption except those stated in its own text or in an assumption it
cites by number.

Every result is numbered within its section, so that Theorem~7.3.2 is the second
numbered statement of Section~7.3. Cross-references are hyperlinked. The two
indices at the end of the volume list, respectively, defined terms and named
results, and the symbols with the page on which each is introduced.

\chapter*{Mathematical Prerequisites}
\addcontentsline{toc}{chapter}{Mathematical Prerequisites}
\markboth{Mathematical Prerequisites}{Mathematical Prerequisites}

The volume is self-contained in the sense that every object it uses is defined
before it is used. It is not self-contained in the sense of assuming no prior
exposure: a reader meeting the derivative or the determinant for the first time
here will find the treatment too fast. What is assumed is the content of a solid
undergraduate program in mathematics for economics, namely

\begin{enumerate}[(i)]
	\item elementary set notation and the manipulation of quantifiers;
	\item single-variable differential and integral calculus at the computational
	      level --- differentiating and integrating elementary functions, and the
	      statement of the fundamental theorem of calculus;
	\item matrix algebra: multiplication, inversion, determinants, and the solution
	      of linear systems;
	\item elementary probability: random variables, expectation, variance, and the
	      common families of distributions.
\end{enumerate}

Each of these is revisited from a higher standpoint rather than assumed silently.
\autoref{ch:sets} restates set theory in the language of relations and mappings
and constructs the real numbers axiomatically; \autoref{ch:linear} replaces
matrix algebra by the theory of linear maps between vector spaces, recovering
matrices as coordinate representations; \autoref{ch:calculus1} rebuilds the
derivative as a local linear approximation; and \autoref{ch:probability} builds
probability on a measure space. A reader who is comfortable with (i)--(iv) can
read the corresponding chapters quickly, but should not skip them: the higher
standpoint is what the later chapters use.

No prior exposure to metric-space topology, convex analysis, measure theory, or
dynamic programming is assumed.

\chapter*{Notation and Conventions}
\addcontentsline{toc}{chapter}{Notation and Conventions}
\markboth{Notation and Conventions}{Notation and Conventions}

The notation follows the Gerzensee course materials wherever they fix a
convention, and standard mathematical usage elsewhere. Conflicts between sources
have been resolved in favour of the version that makes hypotheses most visible;
where a competing convention is common in economics, it is noted at the point of
first use.

\section*{Sets and numbers}

\begin{center}
	\renewcommand{\arraystretch}{1.25}
	\begin{tabular}{@{}ll@{}}
		\toprule
		Symbol                                       & Meaning                                                        \\
		\midrule
		$\N$, $\Z$, $\Q$, $\R$, $\C$                 & natural numbers (excluding $0$), integers, rationals, reals,
		complex numbers                                                                                               \\
		$\N_0$                                       & $\N\cup\{0\}$                                                  \\
		$\R_+$, $\R_{++}$                            & $[0,\infty)$ and $(0,\infty)$                                  \\
		$\overline{\R}$                              & extended real line $\R\cup\{-\infty,+\infty\}$                 \\
		$\R^n$                                       & $n$-fold Cartesian product of $\R$; elements are column vectors \\
		$\R^n_+$, $\R^n_{++}$                        & vectors with all coordinates $\geq 0$, respectively $>0$        \\
		$\comp{A}$, $A\setminus B$                   & complement of $A$, set difference                              \\
		$\pow{X}$                                    & power set of $X$                                               \\
		$\Int A$, $\closure{A}$, $\partial A$        & interior, closure, boundary of $A$                             \\
		$\relint A$, $\conv A$, $\aff A$             & relative interior, convex hull, affine hull of $A$             \\
		$\one_A$                                     & indicator function of the set $A$                              \\
		$\eqdef$                                     & equality that defines the left-hand side                       \\
		\bottomrule
	\end{tabular}
\end{center}

\section*{Vectors, matrices and maps}

Vectors in $\R^n$ are columns; $x^{\transp}$ denotes the transpose, so that
$x^{\transp}y=\ip{x}{y}=\sum_{i=1}^n x_iy_i$. Matrices are written $A$, $B$,
$X$, with $A^{\transp}$ the transpose, $A^{-1}$ the inverse when it exists,
$\rank A$ the rank, $\tr A$ the trace, $\det A$ the determinant, $\Ker A$ the
kernel and $\Image A$ the image. The $i$-th standard basis vector of $\R^n$ is
$e^i$. For a matrix $A$, the notation $A\succeq 0$ ($A\succ 0$) means that $A$
is symmetric positive semidefinite (positive definite).

For $f\colon A\to B$ the notation $f(S)$ denotes the image of $S\subseteq A$ and
$f^{-1}(T)$ the preimage of $T\subseteq B$; $f^{-1}$ is used for the inverse map
only where $f$ is stated to be a bijection. A correspondence, that is a
set-valued map, is written $\Gamma\colon A\rightrightarrows B$ and its graph is
$\gph\Gamma$.

Derivatives use $f'(x)$ in one variable, $\nabla f(x)\in\R^n$ for the gradient of
a real-valued $f$ on $\R^n$, $Df(x)$ for the $m\times n$ Jacobian matrix of
$f\colon\R^n\to\R^m$, and $\Hess f(x)$ for the $n\times n$ Hessian. The gradient
is a column vector and $Df(x)=\nabla f(x)^{\transp}$ when $m=1$.

\section*{Analysis}

A metric space is a pair $(X,d)$ and a normed space a pair $(X,\norm{\cdot})$;
the pair is abbreviated to $X$ when the metric or norm is clear. Open balls are
$\ball{\varepsilon}{x}=\{y\in X: d(x,y)<\varepsilon\}$. Sequences are
$(x_n)_{n\in\N}$ and subsequences $(x_{n_k})_{k\in\N}$. The spaces of continuous
and of bounded functions from $X$ to $Y$ are $\Cspace{X}{Y}$ and
$\Bspace{X}{Y}$; both carry the supremum norm
$\norm{f}_\infty=\sup_{x\in X}\abs{f(x)}$ unless stated otherwise.

\section*{Optimisation}

A constrained problem is written
\[
	\max_{x\in\R^m} f(x)\quad \st\quad g_j(x)\geq 0,\ j=1,\dots,n,
\]
with the constraints in the $\geq 0$ orientation used throughout the Gerzensee
course \autocite[slide 3]{fukuda2026optimization}. The feasible set is $D$, the
Lagrangian is $\Lag$, and multipliers are $y=(y_1,\dots,y_n)\in\R^n_+$. Writing
constraints as $g_j\geq0$ rather than $g_j\leq0$ fixes the sign of the multiplier
at a maximum to be non-negative, which is what makes the shadow-price reading
correct without a case distinction; the opposite convention, common in the
optimisation literature, is noted in \autoref{ch:optimization}.

\section*{Probability}

A probability space is $(\Omega,\mathcal{F},\Prob)$, random variables are
$X,Y,Z$, and realisations are $x,y,z$. Expectation, variance and covariance are
$\E$, $\Var$ and $\Cov$; $X\indep Y$ denotes independence. Conditional
expectation given an event $H$ is $\E[X\mid H]$ and given a random variable $Y$
is $\E[X\mid Y]$, the latter being a random variable, that is, a function of $Y$.
Precision is $\tau=1/\sigma^2$, following the treatment of normal-normal updating
in \autocite[slides 27--28]{fukuda2026probability}.

\section*{Dynamics}

Time is $t\in\N_0$ in discrete time and $t\in[0,\infty)$ in continuous time.
State variables are $x$ and controls $u$; the discount factor is
$\beta\in(0,1)$; the transition is $x'=g(x,u)$, with a prime denoting the
next period, as in \autocite[slide 8]{fukuda2026dynamic}. The value function is
$V$, the policy function is $\policy$, and the Bellman operator is $T$.
